\documentclass[12pt,a4paper]{article}
\usepackage[margin=25mm, headheight=15pt, headsep=10pt]{geometry}
\usepackage{fontspec}
\usepackage[table]{xcolor} % Note: table option is required for rowcolors
\usepackage{titlesec}
\usepackage{tocloft}
\usepackage{fancyhdr}
\usepackage{booktabs}
\usepackage{tcolorbox}
\tcbuselibrary{skins, breakable}
\usepackage[colorlinks=true, linkcolor=emerald, urlcolor=emerald, bookmarks=true]{hyperref}

\definecolor{emerald}{RGB}{0,102,51}
\definecolor{darkred}{RGB}{139,0,0}
\definecolor{lightgray}{RGB}{245,245,245}

\usepackage[nil]{babel}
\babelprovide[import, main]{bengali}
\babelprovide[import]{arabic}
\babelfont{rm}[Renderer=HarfBuzz,Scale=1.0]{Noto Serif Bengali}
\babelfont{sf}[Renderer=HarfBuzz]{Noto Sans Bengali}
\babelfont{tt}[Renderer=HarfBuzz]{Noto Sans Bengali}
\babelfont[arabic]{rm}[Renderer=HarfBuzz,Scale=1.1]{Noto Naskh Arabic}

% Section styling
\titleformat{\section}{\Large\bfseries\color{emerald}}{\thesection.}{0.5em}{}
\titleformat{\subsection}{\large\bfseries\color{emerald!85!black}}{\thesubsection.}{0.5em}{}
\titleformat{\subsubsection}{\normalsize\bfseries\color{emerald!70!black}}{\thesubsubsection.}{0.5em}{}
\titlespacing*{\section}{0pt}{18pt}{8pt}

% Fancy Header/Footer setup
\pagestyle{fancy}
\fancyhf{} % clear all header and footer fields
\fancyhead[L]{\color{emerald}\small\textbf{\nouppercase{\leftmark}}}
\fancyfoot[L]{\scriptsize\color{black!50}{\lat{AI}}-সংকলিত, আলেম-যাচাইকৃত নয়}
\fancyfoot[C]{\color{emerald!70!black}\thepage}
\renewcommand{\headrulewidth}{0.4pt}
\renewcommand{\headrule}{\hbox to\headwidth{\color{emerald}\leaders\hrule height \headrulewidth\hfill}}
\renewcommand{\footrulewidth}{0pt}

% Custom boxes
% 1. Ayat/Hadith Box
\newtcolorbox{ayatbox}[1][]{
    enhanced, breakable, 
    colback=emerald!5, 
    colframe=emerald, 
    boxrule=0.5pt, 
    arc=3pt, 
    left=10pt, right=10pt, top=10pt, bottom=10pt,
    #1
}

% 2. Quote/Translation Box
\newtcolorbox{quotebox}[1][]{
    enhanced, breakable,
    frame hidden,
    colback=lightgray,
    borderline west={3pt}{0pt}{emerald},
    left=15pt, right=10pt, top=10pt, bottom=10pt,
    sharp corners,
    #1
}

% 3. AI-generated notice box
\newtcolorbox{warnbox}[1][]{
    enhanced, breakable,
    colback=red!4,
    colframe=darkred,
    boxrule=0.8pt,
    arc=3pt,
    left=10pt, right=10pt, top=10pt, bottom=10pt,
    #1
}

% Commands
\newcommand{\arab}[1]{\foreignlanguage{arabic}{#1}}
\newcommand{\kib}[1]{\textcolor{emerald}{\textbf{#1}}}

% Latin (English) fallback: Noto Serif Bengali has NO Latin glyphs
\newfontfamily\latfont[Renderer=HarfBuzz]{Noto Serif}
\newcommand{\lat}[1]{{\latfont #1}}

\setlength{\parskip}{5pt}
\setlength{\parindent}{0pt}

% Fix for missing bullet in Noto Serif Bengali
\renewcommand{\labelitemi}{$\bullet$}



\begin{document}
\begin{center}{\Large\bfseries\color{emerald} \lat{01}-সূরা-ফাতিহা-শব্দ-অর্থ}\end{center}
\vspace{1em}
\section*{সূরা আল-ফাতিহা — শব্দ-শব্দ অর্থ ও ব্যাখ্যা (নামাজের রুকন)}

\begin{warnbox}
 \textbf{গুরুত্বপূর্ণ নোটিশ (\lat{Important} \lat{Notice}):} এই কন্টেন্টটি \lat{AI} (কৃত্রিম বুদ্ধিমত্তা) দিয়ে প্রস্তুত ও সংকলিত। \lat{AI} কঠোর সূত্র-মান নীতি মেনে শুধু নির্ভরযোগ্য উৎস থেকে তথ্য সংগ্রহ করেছে — কেবল সহীহ/হাসান হাদিস ও সহীহ সনদের আসার রাখা হয়েছে, দুর্বল সনদ বাদ দেওয়া হয়েছে। তবে এটি কোনো আলেম বা মুহাদ্দিস কর্তৃক যাচাইকৃত নয়।
\end{warnbox}

\medskip\hrule\medskip

\subsection*{১. সূত্র কার্ড}

\begin{table}[h]\centering\small
\begin{tabular}{ll}
বিষয় & বিবরণ \\
\hline
\textbf{সূরা} & আল-ফাতিহা (১) — মক্কী \\
\textbf{আয়াত সংখ্যা} & ৭ আয়াত \\
\textbf{পারা} & ১ \\
\textbf{নামের অর্থ} & ফাতিহা = উন্মুখ/প্রস্তাবনা — কুরআনের দ্বার;\textit{\lat{salah}} (নামাজ)-এর \textit{\lat{pillar}} (রুকন) \\
\textbf{বৈশিষ্ট্য} & \textbf{\lat{“}যে ব্যক্তি সূরা ফাতিহা পড়েনি, তার নামাজ নেই\lat{”}} — সহীহ বুখারি ৭৫৬, সহীহ মুসলিম ৩৯৪ \\
\textbf{ব\lat{astanza} নাম} & উম্মুল কিতাব (কিতাবের মাতা), সাবা'ুল মഷানি (সাতটি ঘুরে পাওয়ার আয়াত), আল-হাম্দ (প্রশংসা), অশ-শিফা (চিকিৎসা) \\
\hline
\end{tabular}
\end{table}

\medskip\hrule\medskip

\subsection*{২. প্রসঙ্গ ও প্রেক্ষাপট}

সূরা ফাতিহা কুরআনের \textbf{প্রথম সূরা} এবং নামাজের প্রতিটি রাকাতে পঠন \textbf{বাধ্যতামূলক (ফরজ/রুকন)}। নবীজি (সা.) ফরমালেন: \textbf{\lat{“}নামাজ পড়ো, যেমন আমাকে নামাজ পড়তে দেখেছ\lat{”}} (সহীহ বুখারি ৬৩১) — এই আদেশের কোথায় এই সূরাটি প্রতিটি রাকাতে আবর্তন করা হয়। এটি বান্দার আকাঙ্ক্ষা ও আল্লাহর প্রতিশ্রুতির \textbf{\lat{dialogue} (সদ্দша)} — বান্দা माँগে, আল্লাহ দেন।

\medskip\hrule\medskip

\subsection*{৩. সম্পূর্ণ আরবি টেক্সট (মূল)}

\begin{center}\arab{بِسْمِ اللَّهِ الرَّحْمَٰنِ الرَّحِيمِ}\end{center}
\begin{center}\arab{الْحَمْدُ لِلَّهِ رَبِّ الْعَالَمِينَ}\end{center}
\begin{center}\arab{الرَّحْمَٰنِ الرَّحِيمِ}\end{center}
\begin{center}\arab{مَالِكِ يَوْمِ الدِّينِ}\end{center}
\begin{center}\arab{إِيَّاكَ نَعْبُدُ وَإِيَّاكَ نَسْتَعِينُ}\end{center}
\begin{center}\arab{اهْدِنَا الصِّرَاطَ الْمُسْتَقِيمَ}\end{center}
\begin{center}\arab{صِرَاطَ الَّذِينَ أَنْعَمْتَ عَلَيْهِمْ غَيْرِ الْمَغْضُوبِ عَلَيْهِمْ وَلَا الضَّالِّينَ}\end{center}

\medskip\hrule\medskip

\subsection*{৪. বাংলা উচ্চারণ (\lat{Transliteration})}

\textbf{বিসমিল্লাহির রাহমানির রাহীম}
\textbf{আলহামদু লিল্লাহি রাব্বিল আলামীন}
\textbf{আর রাহমানির রাহীম}
\textbf{মালিকি যাওমিদ দীন}
\textbf{ইয়াকা নাআবুদু ওয়া ইয়াকা নাস্তাঈন}
\textbf{ইহদিনাস সিরাতাল মুস্তাকীম}
\textbf{সিরাতাল্লাজীনা আনআমতা আলাইহিম গৈরিল মাগদুবি আলাইহিম ওয়ালাদ্বাল্লীন}

\medskip\hrule\medskip

\subsection*{৫. শব্দ-শব্দ অর্থ (\lat{Word-by-Word} \lat{Vocabulary} \lat{Table})}

\begin{table}[h]\centering\small
\begin{tabular}{lll}
আরবি শব্দ & বাংলা অর্থ & \lat{English} (\lat{learning}) \\
\hline
\textbf{\arab{بِسْمِ}} & নামে & \textit{\lat{in} \lat{the} \lat{name}} \\
\textbf{\arab{اللَّهِ}} & আল্লাহর & \textit{\lat{Allah}} \\
\textbf{\arab{الرَّحْمَٰنِ}} & পরম করুণাময় & \textit{\lat{the} \lat{Most} \lat{Gracious}} \\
\textbf{\arab{الرَّحِيمِ}} & পরম দয়ালু & \textit{\lat{the} \lat{Most} \lat{Merciful}} \\
\textbf{\arab{الْحَمْدُ}} & সকল প্রশংসা & \textit{\lat{all} \lat{praise}} \\
\textbf{\arab{لِلَّهِ}} & আল্লাহর জন্যে & \textit{\lat{for} \lat{Allah}} \\
\textbf{\arab{رَبِّ}} & পালনকর্তা/রব & \textit{\lat{Lord}, \lat{Sustainer}} \\
\textbf{\arab{الْعَالَمِينَ}} & সকল जगতের (জগতসমূহ) & \textit{\lat{the} \lat{worlds}} \\
\textbf{\arab{الرَّحْمَٰنِ}} & পরম করুণাময় & \textit{\lat{the} \lat{Most} \lat{Gracious}} \\
\textbf{\arab{الرَّحِيمِ}} & পরম দয়ালু & \textit{\lat{the} \lat{Most} \lat{Merciful}} \\
\textbf{\arab{مَالِكِ}} & মালিক/স্বামী & \textit{\lat{Master}, \lat{Owner}} \\
\textbf{\arab{يَوْمِ}} & দিন & \textit{\lat{Day}} \\
\textbf{\arab{الدِّينِ}} & বিচার/দীন & \textit{\lat{Judgement}/\lat{Religion}} \\
\textbf{\arab{إِيَّاكَ}} & শুধু তোমাকেই & \textit{\lat{You} \lat{alone}} \\
\textbf{\arab{نَعْبُدُ}} & আমরা ইবাদত করি & \textit{\lat{we} \lat{worship}} \\
\textbf{\arab{وَ}} & ও & \textit{\lat{and}} \\
\textbf{\arab{إِيَّاكَ}} & শুধু তোমাকেই & \textit{\lat{You} \lat{alone}} \\
\textbf{\arab{نَسْتَعِينُ}} & আমরা সাহায্য চাই & \textit{\lat{we} \lat{seek} \lat{help}} \\
\textbf{\arab{اهْدِنَا}} & হেদায়াত দাও/পথ দেখাও & \textit{\lat{guide} \lat{us}} \\
\textbf{\arab{الصِّرَاطَ}} & পথটি & \textit{\lat{the} \lat{path}} \\
\textbf{\arab{الْمُسْتَقِيمَ}} & সরল/সীधা & \textit{\lat{straight}} \\
\textbf{\arab{صِرَاطَ}} & পথ & \textit{\lat{path}} \\
\textbf{\arab{الَّذِينَ}} & যারা & \textit{\lat{those} \lat{who}} \\
\textbf{\arab{أَنْعَمْتَ}} & তুমি অনুগ্রহ করেছ & \textit{\lat{You} \lat{bestowed} \lat{favour}} \\
\textbf{\arab{عَلَيْهِمْ}} & তাদের ওপর & \textit{\lat{upon} \lat{them}} \\
\textbf{\arab{غَيْرِ}} & না/বর্জিত & \textit{\lat{not}} \\
\textbf{\arab{الْمَغْضُوبِ}} & কূপিত/রোষিত & \textit{\lat{those} \lat{who} \lat{earned} \lat{anger}} \\
\textbf{\arab{عَلَيْهِمْ}} & তাদের ওপর & \textit{\lat{upon} \lat{them}} \\
\textbf{\arab{وَلَا}} & এবং না & \textit{\lat{nor}} \\
\textbf{\arab{الضَّالِّ}ิน\arab{َ}} & পথভ্রষ্টরা & \textit{\lat{those} \lat{who} \lat{go} \lat{astray}} \\
\hline
\end{tabular}
\end{table}

\medskip\hrule\medskip

\subsection*{৬. আয়াত-वार পূর্ণ বাংলা অনুবাদ (\lat{Pure} \lat{Bengali})}

\textbf{আয়াত ১:} \textbf{\lat{“}আল্লাহর নামে, যিনি পরম করুণাময়, পরম দয়ালু।\lat{”}}

\textbf{আয়াত ২:} \textbf{\lat{“}সকল প্রশংসা আল্লাহর জন্যে, যিনি সকল জগতের পালনকর্তা।\lat{”}}

\textbf{আয়াত ৩:} \textbf{\lat{“}পরম করুণাময়, পরম দয়ালু।\lat{”}}

\textbf{আয়াত ৪:} \textbf{\lat{“}বিচারের দিনের মালিক।\lat{”}}

\textbf{আয়াত ৫:} \textbf{\lat{“}শুধু তোমাকেই আমরা ইবাদত করি এবং শুধু তোমাকেই আমরা সাহায্য চাই।\lat{”}}

\textbf{আয়াত ৬:} \textbf{\lat{“}আমাদের সরল পথে হেদায়াত কর।\lat{”}}

\textbf{আয়াত ৭:} \textbf{\lat{“}যারা তোমার অনুগ্রহ পেয়েছেন তাদের পথ, তাদের নয় যারা কোপিত हुए, এবং 나 পথভ্রষ্টরাও নয়।\lat{”}}

\medskip\hrule\medskip

\subsection*{৭. সংক্ষিপ্ত ব্যাখ্যা (\lat{Khushu-focused} — \lat{with} \lat{English} \lat{Learning} \lat{Mix})}

\subsubsection*{\textbf{আয়াত ১-২: হামদ (প্রশংসা) — \textit{\lat{Gratitude} (কৃতজ্ঞতা)} শুরু}}
নামাজ শুরু করছি আল্লাহর প্রশংসা দিয়ে। \textbf{\lat{Rabb} (রব)} মানে যিনি সৃষ্টি করলেন, পোষণ করলেন, \lat{dirigida} করেন। \textit{"\lat{Al-alamin}"} — সকল জগত (মানুষ, জিন, ফেরেশতা, প্রাণী, গ্রহ-নক্ষত্র)। এখানে বান্দা \lat{acept}ó করছে: আমার যা কিছু আছে, সব তারই দেওয়া।

\subsubsection*{\textbf{আয়াত ৩: রাহমান ও রাহীম — \textit{\lat{Mercy} (দয়া)}-র দু 형\lat{oyle}}}
\textbf{\lat{Ar-Rahman}} — সকল সৃষ্টির জন্য সাধারণ দয়া (দুনিয়া-আখেরাত)। \textbf{\lat{Ar-Rahim}} — মুমিনদের জন্য বিশেষ দয়া (আখেরাতে)। দুটোর মিল নামাজে \textit{\lat{hope} (আশা)} জাগায়।

\subsubsection*{\textbf{আয়াত ৪: মালিকি যাওমিদ দীন — \textit{\lat{Accountability} (দায়বদ্ধতা)} স্মরণ}}
"\lat{Malik}" — মালিক, বিচারক। "\lat{Yawmid} \lat{Din}" — বিচারের দিন। এই আয়াতে মানব মনে \textit{\lat{fear} (ভয়)} ও \textit{\lat{responsibility} (দায়িত্ব)} জাগে। নামাজ কেবল ইবাদত নয়, বিচারের প্রস্তুতি।

\subsubsection*{\textbf{আয়াত ৫: ইয়াকা নাআবুদু ওয়া ইয়াকা নাস্তাঈন — \textit{\lat{Tawheed} (তাওহীদ)}-এর সার}}
\textbf{"শুধু তোমাকেই ইবাদত করি"} — ইখলাস (বিশুদ্ধতা)। \textbf{"শুধু তোমাকেই সাহায্য চাই"} — তাওয়াক্কুল (আল্লাহর উপর ভরসা)। এই আয়াত দোটি \textit{\lat{worship} (উপাসনা)} ও \textit{\lat{dependence} (নির্ভরতা)}-র ঘোষণা। নামাজে এই আয়াত পড়ে \textit{\lat{sincerity} (সততা)} নতুন করে যাচাই করুন।

\subsubsection*{\textbf{আয়াত ৬: ইহদিনাস সিরাতাল মুস্তাকীম — \textit{\lat{Guidance} (হেদায়াত)} প্রার্থনা}}
"\lat{Sirat}" — পথ। "\lat{Mustaqim}" — সীধা, বাঁকা নয়। বান্দা প্রতিটি রাকাতে এই দোয়া করে কারণ \textit{\lat{guidance} (হেদায়াত)} সম 무언া নন — প্রতিটি মুহূর্তে প্রয়োজন। এটি সেরা \textbf{\lat{du'a} (দোয়া)} যা আল্লাহ নিজে শিখিয়েছেন।

\subsubsection*{\textbf{আয়াত ৭: সিরাতাল্লাজীনা... — \textit{\lat{Discrimination} (বিভেদ)} জানা}}
তিন শ্রেণী মানুষ:
\begin{enumerate}
\item \textbf{আনামতা আলাইহিম} — যাদের ওপর অনুগ্রহ (নবীগণ, সিদ্ধীকিন, শহীদগণ, সালেহিন — নিসা ৪:৬৯)।
\item \textbf{মাগদুবি আলাইহিম} — যারা জানে কিন্তু মান\lat{ena} (যেমন: যাহুদিারা যों জানি) — \textit{\lat{knowledge} \lat{without} \lat{action}}।
\item \textbf{াদ্দাল্লীন} — যারা জানে না এবং পথভ্রষ্ট (যেমন: নেসারারা যারা জ্ঞান ছাড়া ভ্রান্তে) — \textit{\lat{ignorance}}।
\end{enumerate}
বান্দা প্রার্থনা করছে: আমাকে প্রথম গোষ্ঠীর সাথে রাখ, দ্বিতীয় ও তৃতীয় থেকে বাঁচা।

\medskip\hrule\medskip

\subsection*{৮. খুশুর জন্য মূল পয়েন্টস (\lat{Key} \lat{Points} \lat{for} \lat{Khushu})}

\begin{table}[h]\centering\small
\begin{tabular}{ll}
পয়েন্ট & কীভাবে প্রয়োগ করবেন \\
\hline
\textbf{বিসমিল্লাহ} & প্রতিটি রাকাতে কিছুক্ষণ বিরতি দিয়ে পড়ুন — \textit{"আমি আল্লাহর নামে শুরু করছি"} অনুভব করুন। \\
\textbf{আলহামদু} & হৃদয়ে \textit{\lat{gratitude} (কৃতজ্ঞতা)} আনুন — আজো অঙ্গ-প্রত্যঙ্গ কাজ করছে, তা আল্লাহর দান। \\
\textbf{রাব্বিল আলামীন} & সৃষ্টির বিস্তার চিন্তা করুন — আকাশ, পৃথিবী, সমুদ্র, আপনি যেন ছোট একটি কণিকা। \\
\textbf{মালিকি যাওমিদ দীন} & কিয়ামতের দিন দাঁড়ানোর \textit{\lat{visualization} (কল্পনা)} করুন — পাপ-পুণ্যের হিসাব। \\
\textbf{ইয়াকা নাআবুদু} & হাত বাঁধার \arab{الوقت}ে মনে করুন: \textbf{"আমার কোনো ś\lat{w} \lat{drugi} নেই, শুধু আল্লাহ।"} \\
\textbf{ইহদিনা} & দোয়ার সरे চোখ বন্ধ করুন, হৃদয়ে বলুন: \textbf{"হে আল্লাহ, আমাকে পথ থেকে ভটকে যাওয়া থেকে বাঁচাও।"} \\
\textbf{আমিন} & ফাতিহা শেষে নীরবে \textbf{"আমিন"} বলুন — অর্থ: \textbf{"হে আল্লাহ, কবুল করুন।"} ফেরেশতাও আমিন বলে (বুখারি ৭৮০)। \\
\hline
\end{tabular}
\end{table}

\medskip\hrule\medskip

\subsection*{৯. সংশ্লিষ্ট সহীহ হাদিস ও আয়াত}

\begin{table}[h]\centering\small
\begin{tabular}{lll}
বিষয় & হাদিস/আয়াত & গ্রন্থ ও নম্বর \\
\hline
ফাতিহা ছাড়া নামাজ নেই & \textbf{\lat{“}যে ব্যক্তি সূরা ফাতিহা পড়েনি, তার নামাজ নেই (পূর্ণ হয় না)।\lat{”}} & সহীহ বুখারি ৭৫৬, সহীহ মুসলিম ৩৯৪ \\
ফাতিহা কুরআনের সেরা অংশ & \textbf{\lat{“}আমি তোমাকে সাতটি ঘুরে পাওয়ার আয়াত (সাবা'ুল মশানি) দিলাম যা সে কুরআনের মাতা।\lat{”}} & সহীহ বুখারি ৪৪৭৪ \\
ফাতিহা আধা আল্লাহর জন্য, আধা বান্দার জন্য & হাদিসে কুদসি: \textbf{\lat{“}নামাজ আমি আমার বান্দার সাথে আধা-আধা বাটলাম...\lat{”}} (বান্দা পড়ে "ইয়াকা নাআবুদু" এরপর আল্লাহ বলেন: "এটি আমার বান্দার জন্য এবং তার জন্য যা চায়") & সহীহ মুসলিম ৩৯৫, তিরমিযী ৩১৫০ \\
আমিন ফেরেশতাদের সাথে মিললে গুনাহ মাফ & \textbf{\lat{“}যখন ইমাম 'ওয়ালাদ্বাল্লীন' বলেন, তোমরা আমিন বলো... যার আমিন ফেরেশতাদের সাথে মিলে যাবে, তার আগের গুনাহ মাফ।\lat{”}} & সহীহ বুখারি ৭৮০, सहীহ মুসলিম ৪১০ \\
সুরা ফাতিহা শিফা (চিকিৎসা) & \textbf{\lat{“}সূরা ফাতিহা সব রোগের চিকিৎসা।\lat{”}} & সহীহ বুখারি ৫৭৩৬, সুনানে আবু দাউদ ৩৮৯৪ \\
সিরাতাল মুস্তাকীম = ইসলাম & \textbf{\lat{“}সিরাতাল মুস্তাকীম হলো ইসলাম।\lat{”}} & সুনানে ইবনে মাজাহ ৪২, মুসনাদ আহমদ (সহীহ) \\
আনামতা আলাইহিম = নবী, সিদ্দিক, শহিদ, সালেহ & আয়াত: \textbf{\lat{“}যারা আল্লাহ ও তাঁর রাসূলের আজ্ঞা মেনে চলেন, তারা তাদের সাথে যাঁর সাথে আল্লাহ অনুগ্রহ করেছেন — নবীগণ, সিদ্ধীকিন, শহীদগণ ও সালেহিন।\lat{”}} & সূরা আন-নিসা ৪:৬৯ \\
\hline
\end{tabular}
\end{table}

\medskip\hrule\medskip

\subsection*{১০. অভ্যাসের পরামর্শ (\lat{Practice} \lat{Tips})}

\begin{enumerate}
\item \textbf{প্রতিদিন ১ রাকাত \lat{vaqt}} — শুধু ফাতিহা ধীরে ধীরে অর্থ মনে করেই পড়ুন (৫ মিনিট)।
\item \textbf{শব্দ টেবিল মেমোরাইজ করুন} — প্রতিটি আরবি শব্দের অর্থ জানলে নামাজে \textit{\lat{concentration} (মনোযোগ)} বাড়ে।
\item \textbf{রেকর্ড করুন} — নিজের কণ্ঠে ফাতিহা রেকর্ড করে শোনুন, উচ্চারণ শুদ্ধ করুন।
\item \textbf{নামাজে প্রতিটি আয়াতের পর ২-৩ সেকেন্ড বিরতি} — অর্থ মনে প\lat{es}í\lat{a} করার জন্য (ইতমিনান)।
\item \textbf{সিজদায় এই দোয়া করুন}: \textit{"হে আল্লাহ, আমাকে ফাতিহার অর্থ বুঝিয়ে নামাজ দাও।"}
\end{enumerate}
\medskip\hrule\medskip

\textit{শেষ আপডেট: আগস্ট ২০২৬}

\end{document}
